% © 2026 Ing. David Jaroš — CC BY-NC-ND 4.0
%
% This work is licensed under a Creative Commons Attribution-NonCommercial-NoDerivatives
% 4.0 International License (CC BY-NC-ND 4.0).
%
% See LICENSE.md for full license text.

\documentclass[12pt,a4paper]{article}
\usepackage[a4paper,margin=2.5cm]{geometry}
\usepackage{amsmath,amssymb,tcolorbox}
\tcbuselibrary{skins,breakable}
\newtcolorbox{statusbox}[1][]{colback=gray!8!white,colframe=black!60,arc=3pt,boxrule=1pt,title=Status Summary,#1}
\title{Gap EW-2: $\Theta_0$ VEV as SU(2)$_L$ doublet}
\author{Ing.\ David Jaroš}
\date{2026-05-12}
\begin{document}
\maketitle

From the canonical action of SU(2)$_L$ on $\Theta\in\mathrm{Mat}(2,\mathbb C)$,
a pure identity ansatz $\Theta_0\propto\mathbf 1_2$ transforms as singlet, not as
an EW Higgs-like doublet.

A column-projection ansatz
\[
\Theta_0 = (\phi,0),\qquad \phi\in\mathbb C^2,
\]
can realise a doublet-like transformation pattern, but closure to a fully derived
SSB theorem from $S[\Theta]$ is not yet established in canonical files.

Mass relations remain conditional:
\[
M_W = \frac{gv}{2},\qquad M_Z = \frac{M_W}{\cos\theta_W},
\]
with $g,v$ not first-principles derived here.

\begin{statusbox}
\textbf{Hlavní výsledek}: Identitní VEV je singlet; doubletová projekce je kandidát, nikoli uzavřený důkaz z $S[\Theta]$.\\
\textbf{Klasifikace}: $\Theta_0$ jako dublet [MC/OPEN]; W/Z hmoty [CONDITIONAL].\\
\textbf{Dopad na teorii}: Gap EW-2 zůstává otevřený, ale je zúžen na konkrétní reprezentační/projekční problém.\\
\textbf{Příští krok}: Odvodit přesnou transformaci a potenciálový minimizační argument pro dubletový sektor přímo ze spektrální/UBT akce.\\
\textbf{Alpha je odvozena}: NE (unconditional)
\end{statusbox}

\end{document}
